\documentclass[a4paper,12pt]{report}

% Packages needed
\usepackage[utf8]{inputenc}
\usepackage[a4paper, margin=2.5cm]{geometry}
\usepackage{graphicx}
\usepackage{float}
\usepackage{booktabs}
\usepackage{listings}
\usepackage{xcolor}
\usepackage{hyperref}
\usepackage{capt-of}
\usepackage{titlesec}
\usepackage{wrapfig}
\usepackage{needspace}

% Reduce the large empty space before an unnumbered chapter title
\titleformat{\chapter}[display]
  {\normalfont\huge\bfseries}{}{0pt}{\Huge}
\titlespacing*{\chapter}{0pt}{0pt}{30pt}

% Code listing style (Python)
\lstset{
  language=Python,
  basicstyle=\ttfamily\small,
  keywordstyle=\color{blue}\bfseries,
  commentstyle=\color{gray}\ttfamily,
  stringstyle=\color{red}\ttfamily,
  numbers=left,
  numberstyle=\tiny\color{gray},
  stepnumber=1,
  numbersep=5pt,
  frame=single,
  breaklines=true,
  breakatwhitespace=false,
  showstringspaces=false,
  tabsize=2,
  captionpos=b
}

\begin{document}

\chapter*{Text-Guided Object Detection}
\addcontentsline{toc}{chapter}{Text-Guided Object Detection}
\label{chap:text-guided-detection}

\section{Overview}
\label{sec:tgod-overview}

This chapter presents the implementation of the proposed text-guided object
detection system. Unlike conventional object detection models that detect
all objects in an image, the proposed architecture localizes only the
object specified by the user's text prompt. The system combines visual
information extracted from the input image with semantic information
generated from the textual prompt to accurately identify the target
object. The implementation consists of dataset preparation, visual feature
extraction, text encoding, multimodal fusion, object localization, and
model training.

\section{Dataset Preparation}
\label{sec:tgod-dataset-preparation}

The proposed model is trained using images annotated in the YOLO format.
To support text-guided object detection, every annotated object is
converted into an independent training sample. Each sample consists of an
input image, an object prompt, and the corresponding normalized bounding
box.

For example, if an image contains a person, a car, and a dog, three
different training samples are generated from the same image. Although the
image remains unchanged, each sample uses a different text prompt and
bounding box. This strategy enables the model to associate textual
descriptions with their corresponding object locations.

Before training, all images are resized to $224 \times 224$ pixels,
normalized to the range $[0,1]$, and converted into tensors. The object
names are also encoded into numerical tokens that can be processed by the
Text Encoder.

\begin{lstlisting}[caption={Creating one training sample for each object}, label={lst:tgod-training-sample}]
prompt = self.tokenizer.decode(class_id)

self.samples.append({
    "image": image_name,
    "prompt": prompt,
    "bbox": [x, y, w, h],
})
\end{lstlisting}

\begin{lstlisting}[caption={Image preprocessing}, label={lst:tgod-image-preprocessing}]
image = cv2.resize(
    image,
    (IMAGE_SIZE, IMAGE_SIZE)
)
image /= 255.0
\end{lstlisting}

\section{Proposed MiniYOLO Architecture}
\label{sec:tgod-architecture}

The proposed MiniYOLO architecture consists of four main components: the
Backbone Network, the Text Encoder, the Cross-Attention Fusion module, and
the Detection Head.

The Backbone extracts hierarchical visual features from the input image,
while the Text Encoder generates a semantic representation of the object
prompt. Both representations are combined through the Cross-Attention
Fusion module, allowing the network to focus only on image regions related
to the requested object. Finally, the Detection Head predicts the bounding
box coordinates and confidence score.

\begin{lstlisting}[caption={Model architecture}, label={lst:tgod-model-architecture}]
self.backbone = Backbone()
self.text_encoder = TextEncoder()
self.fusion = CrossAttentionFusion()
self.head = DetectionHead()
\end{lstlisting}

\section{Backbone Network}
\label{sec:tgod-backbone}

The Backbone is responsible for extracting high-level visual features from
the input image before the detection stage. It consists of five
convolutional stages, where each stage applies a convolution layer
followed by Batch Normalization, LeakyReLU activation, and Max Pooling.
This design enables the network to progressively learn richer feature
representations while reducing the spatial dimensions of the input.

The early convolutional layers capture low-level features such as edges,
corners, and textures, whereas deeper layers learn more complex semantic
information including object parts, shapes, and overall structures. Batch
Normalization stabilizes the learning process by normalizing feature
distributions, while LeakyReLU helps maintain gradient flow and avoids
inactive neurons during training. Max Pooling reduces the feature map size
while preserving the most informative visual patterns, making the
extracted features more robust and computationally efficient.

Starting from an input image of $224 \times 224 \times 3$, the Backbone
gradually increases the number of feature channels from 32 to 512, while
reducing the spatial resolution to $7 \times 7$. The resulting feature map
provides a compact yet informative representation of the image, which is
then forwarded to the Cross-Attention Fusion module.

\begin{lstlisting}[caption={Backbone Convolution Block}, label={lst:tgod-backbone-block}]
self.block = nn.Sequential(
    nn.Conv2d(...),
    nn.BatchNorm2d(...),
    nn.LeakyReLU(0.1)
)
\end{lstlisting}

\section{Text Encoder}
\label{sec:tgod-text-encoder}

The Text Encoder converts the input object prompt into a semantic feature
vector. First, the object token is transformed into a dense embedding
representation using an Embedding layer. The embedded sequence is then
processed by a Bidirectional LSTM to capture contextual information from
both directions.

The extracted features are summarized using Mean Pooling, normalized with
Layer Normalization, and projected into the same feature space as the
image features. This semantic representation is then forwarded to the
fusion module.

\begin{lstlisting}[caption={Text encoding}, label={lst:tgod-text-encoding}]
self.embedding = nn.Embedding(
    VOCAB_SIZE,
    EMBEDDING_DIM
)

self.encoder = nn.LSTM(
    EMBEDDING_DIM,
    HIDDEN_DIM,
    bidirectional=True
)
\end{lstlisting}

\section{Cross-Attention Fusion}
\label{sec:tgod-cross-attention}

The Cross-Attention Fusion module is the core component of the proposed
architecture because it combines visual and textual information into a
unified representation. Instead of relying only on image features, the
detector uses the object prompt to guide the localization process.

Both image and text features are first projected into the same feature
space. The text feature vector is then expanded across the spatial
dimensions of the image feature map. Element-wise multiplication is
performed to compute an attention representation, which is further
processed to generate an attention map highlighting image regions related
to the requested object. The attention map is finally multiplied with the
image features, enhancing relevant regions while suppressing background
information.

\begin{lstlisting}[caption={Cross-Attention Fusion}, label={lst:tgod-cross-attention-fusion}]
image = self.image_projection(image_features)
text = self.text_projection(text_features)

attention = image * text
fused = image * attention
\end{lstlisting}

\section{Detection Head}
\label{sec:tgod-detection-head}

The Detection Head receives the fused feature representation and predicts
the final localization result. Additional convolutional layers refine the
extracted features before applying Global Average Pooling to obtain a
compact feature vector.

The pooled features are then passed through fully connected layers to
predict five values corresponding to the bounding box center coordinates,
width, height, and confidence score. Since the proposed model detects only
the object specified by the text prompt, the Detection Head produces a
single prediction for each input sample.

\begin{lstlisting}[caption={Detection Head}, label={lst:tgod-detection-head}]
self.pool = nn.AdaptiveAvgPool2d(1)

self.regressor = nn.Sequential(
    nn.Linear(256,128),
    nn.Linear(128,64),
    nn.Linear(64,5)
)
\end{lstlisting}

\section{Model Training}
\label{sec:tgod-model-training}

The model is trained using paired image-text samples generated during
dataset preparation. During each iteration, the input image and its
corresponding text prompt are processed simultaneously to predict the
target bounding box.

Training is performed using the AdamW optimizer together with a learning
rate scheduler. The loss function combines Smooth L1 Loss for bounding box
regression and Binary Cross-Entropy Loss for confidence prediction. The
model with the lowest validation loss is automatically stored as the best
checkpoint and later used during inference.

\begin{lstlisting}[caption={Training step}, label={lst:tgod-training-step}]
predictions = model(images, prompts)

loss, metrics = criterion(
    predictions,
    targets
)

loss.backward()
optimizer.step()
\end{lstlisting}

\end{document}
